%%%%%%%%%%%%%%%%%%%%%%%%%%%%%%%%%
%
%       EXERCÍCIO
%
%%%%%%%%%%%%%%%%%%%%%%%%%%%%%%%%%

\definevalorquestao

\begin{exercicioBanco}[\valorquestao]
O valor médio de comercialização da saca de milho de 60 quilos na Bolsa de Cereais é apresentado abaixo, em reais, para os últimos 40 meses.
%
\begin{center}
    \begin{tabular}{|C{30pt}|C{30pt}|C{30pt}|C{30pt}|C{30pt}|C{30pt}|C{30pt}|C{30pt}|C{30pt}|C{30pt}|}
        \hline
        6.1 & 6.2 & 6.3 & 6.5 & 6.7 & 6.9 & 7.2 & 7.2 & 7.3 & 7.3 \\
        \hline
        7.3 & 7.4 & 7.4 & 7.4 & 7.4 & 7.4 & 7.4 & 7.4 & 7.5 & 7.5 \\
        \hline
        7.5 & 7.5 & 7.5 & 7.5 & 7.6 & 7.6 & 7.6 & 7.6 & 7.6 & 7.6 \\
        \hline
        7.7 & 7.7 & 7.7 & 7.8 & 7.9 & 8.1 & 8.1 & 8.1 & 8.2 & 8.3 \\
        \hline
    \end{tabular}
\end{center}
%
%
\begin{enumerate}[(a)]
    \item Qual é a variável estatística em questão?
    %
    \item Construa o histograma de tal variável estatística, considerando o comprimento de cada faixa igual a \(0.4\), com a primeira faixa começando em \(6\).
    %
    \item Construa o boxplot de tal variável estatística.
\end{enumerate}
%
\end{exercicioBanco}

